\documentclass[12pt,a4paper]{report}
\usepackage[utf8]{inputenc}
\usepackage[margin=1in]{geometry}
\usepackage{xcolor}
\usepackage{listings}
\usepackage{hyperref}
\usepackage{booktabs}
\usepackage{graphicx}
\usepackage{titlesec}
\usepackage{indentfirst}

% Color definitions for code styling
\definecolor{codegreen}{rgb}{0,0.6,0}
\definecolor{codegray}{rgb}{0.5,0.5,0.5}
\definecolor{codepurple}{rgb}{0.58,0,0.82}
\definecolor{backcolour}{rgb}{0.95,0.95,0.96}

% Listings settings
\lstdefinestyle{mystyle}{
    backgroundcolor=\color{backcolour},   
    commentstyle=\color{codegreen},
    keywordstyle=\color{magenta},
    numberstyle=\tiny\color{codegray},
    stringstyle=\color{codepurple},
    basicstyle=\ttfamily\footnotesize,
    breakatwhitespace=false,         
    breaklines=true,                 
    captionpos=b,                    
    keepspaces=true,                 
    numbers=left,                    
    numbersep=5pt,                  
    showspaces=false,                
    showstringspaces=false,
    showtabs=false,                  
    tabsize=4
}
\lstset{style=mystyle}

% Hyperref setup
\hypersetup{
    colorlinks=true,
    linkcolor=blue,
    filecolor=magenta,      
    urlcolor=cyan,
    pdftitle={Django E-Commerce Backend Design Documentation},
    pdfpagemode=FullScreen,
}

% Title details
\title{
    \textbf{Django E-Commerce Backend Design}\\
    \large{Comprehensive Project & Implementation Documentation}\\
    \vspace{1cm}
    \includegraphics[width=0.4\textwidth]{https://via.placeholder.com/150} % Dummy logo/placeholder
}
\author{\textbf{Abdul} \\ Ecommerce Developer}
\date{\today}

\begin{document}

\maketitle

\tableofcontents
\newpage

\chapter{Introduction \& Architecture}

\section{Project Overview}
This project is a complete, progressive implementation of a lightweight E-Commerce Backend application designed and built using the Python Django web framework. The development cycle spans three distinct stages (referred to as Weeks), where the system evolves from a static HTML mock-up using in-memory data to a fully dynamic database-backed web application complete with search systems, pagination, and user authentication management.

\section{Target Audience \& Goals}
The main purpose of this implementation is to showcase:
\begin{itemize}
    \item Progressive feature addition in modern web applications.
    \item Clean separation of concerns using the Model-View-Template (MVT) architecture.
    \item Integration of Django's default SQLite database for dynamic records.
    \item Implementation of standard user flows such as user signup, authenticated product additions, search filtering, and paginated navigation.
\end{itemize}

\section{System Architecture (Django MVT)}
The backend operates strictly under the Django Model-View-Template design pattern:
\begin{description}
    \item[Models] Defines the database schema. In our application, the \texttt{Product} model represents catalog items.
    \item[Views] Handles the request/response lifecycle, queries database tables, implements filters, and passes contextual data to templates.
    \item[Templates] Renders raw HTML files integrated with Django template tag logic (e.g., conditional navigation loops and context injection).
\end{description}

---

\chapter{Progressive Development Stages}

\section{Week 1: Static Layout \& Views}
In the first week of development, focus is placed on design layouts and routing basics.
\begin{itemize}
    \item \textbf{Data Storage:} Memory-based list representation (\texttt{SAMPLE\_PRODUCTS}) inside the views.
    \item \textbf{Endpoints:} Homepage displaying featured items, shop category list, and detailed product profile paths.
    \item \textbf{Limitation:} Changes to the product catalog are not persistent because no database backend is active.
\end{itemize}

\section{Week 2: Dynamic Layout \& Database Integration}
The second stage migrates static arrays into the database engine.
\begin{itemize}
    \item \textbf{Database Engine:} SQLite3 is configured. Django models are compiled, and migrations are applied to create tables.
    \item \textbf{Search Capability:} A search bar checks user input against product names and categories using Django's \texttt{Q} queries (\texttt{name\_\_icontains} and \texttt{category\_\_icontains}).
\end{itemize}

\section{Week 3: Final Production Features}
The final week optimizes UX and security.
\begin{itemize}
    \item \textbf{Pagination:} Prevents page bloating by limiting visibility to two products per page. This enforces pagination logic and simplifies tests.
    \item \textbf{Authentication System:} Implements core login/signup components.
    \item \textbf{Authorization Guards:} Adds security rules prohibiting anonymous guests from posting new items via \texttt{@login\_required} decorators on writing views.
\end{itemize}

---

\chapter{Database Schema \& Design}

\section{Product Model Fields}
The primary database entity is the \texttt{Product} table. The attributes are defined in \texttt{models.py} as follows:

\begin{table}[h]
\centering
\begin{tabular}{llll}
\toprule
\textbf{Field Name} & \textbf{Data Type} & \textbf{Constraints} & \textbf{Purpose} \\ \midrule
\texttt{id} & AutoField (Primary Key) & Auto-increment & Unique ID of product \\
\texttt{name} & CharField & Max length = 200 & Title of the catalog item \\
\texttt{price} & DecimalField & Max digits = 10, places = 2 & Item price \\
\texttt{category} & CharField & Max length = 100 & Product category grouping \\
\texttt{image} & URLField & Max length = 500 & Link to item cover image \\
\texttt{description} & TextField & No size limit & Detailed descriptions \\
\texttt{stock} & IntegerField & Default = 10 & Number of units in inventory \\
\bottomrule
\end{tabular}
\caption{Product Model Field Definitions}
\label{tab:product_model}
\end{table}

\section{Database Migrations Workflow}
Whenever changes are made to models, Django propagates them through migrations:
\begin{lstlisting}[language=bash]
# Creates migration instructions
python manage.py makemigrations

# Applies structural changes to database file (sqlite3)
python manage.py migrate
\end{lstlisting}

---

\chapter{Core Features Implementation Details}

\section{Search Query Filtering}
We use Django's \texttt{Q} objects to query the backend database for matches against both product titles and category designations in a case-insensitive search format.
\begin{lstlisting}[language=python]
# View snippet executing the search
products_list = Product.objects.filter(Q(name__icontains=query) | Q(category__icontains=query))
\end{lstlisting}

\section{Testing-oriented Pagination}
For efficient pagination testing, the paginator threshold is set to 2 products per page.
\begin{lstlisting}[language=python]
paginator = Paginator(products_list, 2) 
page_number = request.GET.get('page')
page_obj = paginator.get_page(page_number)
\end{lstlisting}

\section{Authorization Guarding}
Adding a product involves updating the database, which requires a validated session. The view is locked using Django's security middleware:
\begin{lstlisting}[language=python]
@login_required
def add_product_view(request):
    # Form processing logic here...
\end{lstlisting}

---

\chapter{Source Code Reference (Week 3 - Final)}

Here are the complete backend source codes for the final week implementation.

\section{Database Models (\texttt{models.py})}
\begin{lstlisting}[language=python, caption=models.py]
from django.db import models

class Product(models.Model):
    name = models.CharField(max_length=200)
    price = models.DecimalField(max_digits=10, decimal_places=2)
    category = models.CharField(max_length=100)
    image = models.URLField(max_length=500, default="https://via.placeholder.com/150")
    description = models.TextField()
    stock = models.IntegerField(default=10)

    def __str__(self):
        return self.name
\end{lstlisting}

\section{System Views Routing Controller (\texttt{views.py})}
\begin{lstlisting}[language=python, caption=views.py]
from django.shortcuts import render, redirect, get_object_or_404
from django.contrib.auth.forms import UserCreationForm
from django.contrib.auth.decorators import login_required
from django.core.paginator import Paginator
from django.db.models import Q
from .models import Product

# 1. Home View
def home_view(request):
    featured_products = Product.objects.all()[:2]
    return render(request, 'index.html', {'featured_products': featured_products})

# 2. Product List View with Search AND Pagination
def product_list_view(request):
    query = request.GET.get('q', '')
    if query:
        products_list = Product.objects.filter(Q(name__icontains=query) | Q(category__icontains=query))
    else:
        products_list = Product.objects.all()

    # Pagination: 2 products per page
    paginator = Paginator(products_list, 2) 
    page_number = request.GET.get('page')
    page_obj = paginator.get_page(page_number)

    return render(request, 'products.html', {'page_obj': page_obj, 'query': query})

# 3. Product Detail View
def product_detail_view(request, product_id):
    product = get_object_or_404(Product, pk=product_id)
    return render(request, 'product-detail.html', {'product': product})

# 4. Signup View
def signup_view(request):
    if request.method == 'POST':
        form = UserCreationForm(request.POST)
        if form.is_valid():
            form.save()
            return redirect('login')
    else:
        form = UserCreationForm()
    return render(request, 'registration/signup.html', {'form': form})

# 5. Add Product View (Restricted to logged-in users)
@login_required
def add_product_view(request):
    if request.method == 'POST':
        name = request.POST['name']
        price = request.POST['price']
        category = request.POST['category']
        description = request.POST['description']
        stock = request.POST['stock']
        image = request.POST.get('image', '')
        if not image.strip():
            image = "https://via.placeholder.com/150"
        
        Product.objects.create(
            name=name, 
            price=price, 
            category=category, 
            description=description, 
            stock=stock,
            image=image
        )
        return redirect('product_list')
        
    return render(request, 'add-product.html')
\end{lstlisting}

\section{Endpoint URL Mapping Controller (\texttt{urls.py})}
\begin{lstlisting}[language=python, caption=urls.py]
from django.contrib import admin
from django.urls import path, include
from django.contrib.auth import views as auth_views
from . import views

urlpatterns = [
    path('admin/', admin.site.urls),
    path('', views.home_view, name='home'),
    path('products/', views.product_list_view, name='product_list'),
    path('products/<int:product_id>/', views.product_detail_view, name='product_detail'),
    path('products/add/', views.add_product_view, name='add_product'),
    
    # Auth Routes
    path('login/', auth_views.LoginView.as_view(template_name='registration/login.html'), name='login'),
    path('logout/', auth_views.LogoutView.as_view(), name='logout'),
    path('signup/', views.signup_view, name='signup'),
]
\end{lstlisting}

---

\chapter{Installation \& Execution Guide}

\section{Environment Configuration}
To deploy the backend database server locally:
\begin{enumerate}
    \item Clone or download project files.
    \item Install packages needed using the terminal window:
    \begin{lstlisting}[language=bash]
    pip install django
    \end{lstlisting}
    \item Change directory to Week 3 (Final implementation):
    \begin{lstlisting}[language=bash]
    cd week-3-final
    \end{lstlisting}
\end{enumerate}

\section{Database setup and Superuser Creation}
Generate the schema tables and configure custom login profiles:
\begin{lstlisting}[language=bash]
# Compile DB structural mapping 
python manage.py migrate

# Configure Admin account credentials
python manage.py createsuperuser
\end{lstlisting}

\section{Run Development Server}
Run the web application locally:
\begin{lstlisting}[language=bash]
python manage.py runserver
\end{lstlisting}
The local address is \url{http://127.0.0.1:8000/}.

---

\chapter{Conclusion}
This Django project serves as a step-by-step model showing how static websites evolve into dynamic e-commerce solutions. The current version builds a solid base with authentication, product forms, pagination, and data querying. Future versions can expand on this by adding shopping carts, payment gateways (like Stripe), and responsive tailwind UI styles.

\end{document}
